\appendix
\chapter{附录}

\section{附录 A TaskContract 示例}

TaskContract 包含 task\_id、steps、skill、parameters、preconditions、timeout、safety\_constraints 和 version 等字段；所有候选动作必须通过 schema、语义校验和 SafetyShield。\par

\section{附录 B 状态机}

边缘任务状态机覆盖 CREATED、RUNNING、PAUSED、RECOVERING、COMPLETED、FAILED 和 SAFETY\_STOPPED；Simulation Runtime job 状态机覆盖 QUEUED、LEASED、RUNNING、FINALIZING、SUCCEEDED、FAILED、CANCELLED、TIMED\_OUT、INTERRUPTED 和 RECOVERY\_PENDING。\par

\section{附录 C F01-F20 与 B01-B03}

\{\{ experiment\_status \}\}\par

\noindent $\bullet$ B01\_LLM\_ONLY\_ONESHOT：一次模型调用生成完整 TaskContract，异常时失败或整体重试。\par
\noindent $\bullet$ B02\_LLM\_ONLY\_REACTIVE：每一步或异常后调用模型生成下一动作，仍经过契约校验和 SafetyShield。\par
\noindent $\bullet$ B03\_PIPELINE\_ONLY\_PAIRED\_DESIGN：fake-provider 仅用于验证配对对比管线，不进入正式模型性能结论。\par
\noindent $\bullet$ B03\_REAL\_RUNTIME\_PAIRED\_COMPARISON：仅在 REAL\_LLM\_RUNTIME 或 LOCAL\_LLM\_RUNTIME accepted 后，用于 LLM-only one-shot、reactive、PCSC、ETEAC、AUTO 的相同场景、seed、后端和安全策略配对比较。\par

\section{附录 D 图表索引}

\noindent $\bullet$ 图1 system\_architecture：thesis/figures/svg/system\_architecture.svg（source=design\_diagram，authority=L1，formal\_allowed=true）\par
\noindent $\bullet$ 图2 pcsc\_flow：thesis/figures/svg/pcsc\_flow.svg（source=design\_diagram，authority=L1，formal\_allowed=true）\par
\noindent $\bullet$ 图3 eteac\_flow：thesis/figures/svg/eteac\_flow.svg（source=design\_diagram，authority=L1，formal\_allowed=true）\par
\noindent $\bullet$ 图4 auto\_decision：thesis/figures/svg/auto\_decision.svg（source=design\_diagram，authority=L1，formal\_allowed=true）\par
\noindent $\bullet$ 图5 llm\_only\_oneshot：thesis/figures/svg/llm\_only\_oneshot.svg（source=design\_diagram，authority=L1，formal\_allowed=true）\par
\noindent $\bullet$ 图6 llm\_only\_reactive：thesis/figures/svg/llm\_only\_reactive.svg（source=design\_diagram，authority=L1，formal\_allowed=true）\par
\noindent $\bullet$ 图7 safetyshield\_flow：thesis/figures/svg/safetyshield\_flow.svg（source=design\_diagram，authority=L1，formal\_allowed=true）\par
\noindent $\bullet$ 图8 edge\_state\_machine：thesis/figures/svg/edge\_state\_machine.svg（source=design\_diagram，authority=L1，formal\_allowed=true）\par
\noindent $\bullet$ 图9 local\_replanning：thesis/figures/svg/local\_replanning.svg（source=design\_diagram，authority=L1，formal\_allowed=true）\par
\noindent $\bullet$ 图10 experiment\_matrix：thesis/figures/svg/experiment\_matrix.svg（source=design\_diagram，authority=L1，formal\_allowed=true）\par
\noindent $\bullet$ 图11 phase\_roadmap：thesis/figures/svg/phase\_roadmap.svg（source=design\_diagram，authority=L1，formal\_allowed=true）\par
\noindent $\bullet$ 图12 ablation\_comparison：thesis/figures/svg/ablation\_comparison.svg（source=aggregate\_payload，authority=L4，formal\_allowed=true）\par
\noindent $\bullet$ 图13 backend\_delta\_distribution：thesis/figures/svg/backend\_delta\_distribution.svg（source=aggregate\_payload，authority=L4，formal\_allowed=true）\par
\noindent $\bullet$ 图14 cloud\_invocation\_bar：thesis/figures/svg/cloud\_invocation\_bar.svg（source=aggregate\_payload，authority=L4，formal\_allowed=true）\par
\noindent $\bullet$ 图15 communication\_count\_bar：thesis/figures/svg/communication\_count\_bar.svg（source=aggregate\_payload，authority=L4，formal\_allowed=true）\par
\noindent $\bullet$ 图16 completion\_time\_box：thesis/figures/svg/completion\_time\_box.svg（source=aggregate\_payload，authority=L4，formal\_allowed=true）\par
\noindent $\bullet$ 图17 failure\_composition：thesis/figures/svg/failure\_composition.svg（source=aggregate\_payload，authority=L4，formal\_allowed=true）\par
\noindent $\bullet$ 图18 local\_recovery\_effect：thesis/figures/svg/local\_recovery\_effect.svg（source=aggregate\_payload，authority=L4，formal\_allowed=true）\par
\noindent $\bullet$ 图19 mode\_overall\_comparison：thesis/figures/svg/mode\_overall\_comparison.svg（source=aggregate\_payload，authority=L4，formal\_allowed=true）\par
\noindent $\bullet$ 图20 mujoco\_isaac\_paired\_scatter：thesis/figures/svg/mujoco\_isaac\_paired\_scatter.svg（source=aggregate\_payload，authority=L4，formal\_allowed=true）\par
\noindent $\bullet$ 图21 network\_latency\_sensitivity：thesis/figures/svg/network\_latency\_sensitivity.svg（source=aggregate\_payload，authority=L4，formal\_allowed=true）\par
\noindent $\bullet$ 图22 packet\_loss\_sensitivity：thesis/figures/svg/packet\_loss\_sensitivity.svg（source=aggregate\_payload，authority=L4，formal\_allowed=true）\par
\noindent $\bullet$ 图23 planner\_latency：thesis/figures/svg/planner\_latency.svg（source=aggregate\_payload，authority=L4，formal\_allowed=true）\par
\noindent $\bullet$ 图24 planner\_valid\_contract\_rate：thesis/figures/svg/planner\_valid\_contract\_rate.svg（source=aggregate\_payload，authority=L4，formal\_allowed=true）\par
\noindent $\bullet$ 图25 replanning\_count\_comparison：thesis/figures/svg/replanning\_count\_comparison.svg（source=aggregate\_payload，authority=L4，formal\_allowed=true）\par
\noindent $\bullet$ 图26 safety\_intervention\_count：thesis/figures/svg/safety\_intervention\_count.svg（source=aggregate\_payload，authority=L4，formal\_allowed=true）\par
\noindent $\bullet$ 图27 stress\_recovery\_timeline：thesis/figures/svg/stress\_recovery\_timeline.svg（source=aggregate\_payload，authority=L4，formal\_allowed=true）\par
\noindent $\bullet$ 图28 success\_rate\_comparison：thesis/figures/svg/success\_rate\_comparison.svg（source=aggregate\_payload，authority=L4，formal\_allowed=true）\par

\section{附录 E 表格索引}

\noindent $\bullet$ 1. t10\_resource\_overhead：artifacts/phase12\_2\_clean/validation/tables/csv/t10\_resource\_overhead.csv，行数 3\par
\noindent $\bullet$ 2. t11\_failure\_types：artifacts/phase12\_2\_clean/validation/tables/csv/t11\_failure\_types.csv，行数 2\par
\noindent $\bullet$ 3. t12\_reproducibility：artifacts/phase12\_2\_clean/validation/tables/csv/t12\_reproducibility.csv，行数 2\par
\noindent $\bullet$ 4. t1\_system\_capability：artifacts/phase12\_2\_clean/validation/tables/csv/t1\_system\_capability.csv，行数 3\par
\noindent $\bullet$ 5. t2\_mode\_baseline：artifacts/phase12\_2\_clean/validation/tables/csv/t2\_mode\_baseline.csv，行数 3\par
\noindent $\bullet$ 6. t3\_network\_latency：artifacts/phase12\_2\_clean/validation/tables/csv/t3\_network\_latency.csv，行数 3\par
\noindent $\bullet$ 7. t4\_packet\_loss：artifacts/phase12\_2\_clean/validation/tables/csv/t4\_packet\_loss.csv，行数 3\par
\noindent $\bullet$ 8. t5\_fault\_recovery：artifacts/phase12\_2\_clean/validation/tables/csv/t5\_fault\_recovery.csv，行数 3\par
\noindent $\bullet$ 9. t6\_safety\_shield：artifacts/phase12\_2\_clean/validation/tables/csv/t6\_safety\_shield.csv，行数 2\par
\noindent $\bullet$ 10. t7\_backend\_paired：artifacts/phase12\_2\_clean/validation/tables/csv/t7\_backend\_paired.csv，行数 1\par
\noindent $\bullet$ 11. t8\_planner\_provider：artifacts/phase12\_2\_clean/validation/tables/csv/t8\_planner\_provider.csv，行数 3\par
\noindent $\bullet$ 12. t9\_ablation：artifacts/phase12\_2\_clean/validation/tables/csv/t9\_ablation.csv，行数 3\par

\section{附录 F 证据追踪矩阵}

| 第九章 | Phase 12 clean validation 已通过 | status | PHASE12\_VALIDATION\_EXPERIMENTS\_ACCEPTED | artifacts/phase12\_2\_clean/validation/verification/phase12\_summary.json | L4 | validation 级，不代表 full profile。 |\par
| 第九章 | 运行完成数量 | runtime\_completion\_count | 466 | artifacts/phase12\_2\_clean/validation/verification/phase12\_summary.json | L4 | 环境阻塞不计入 runtime completed。 |\par
| 第九章 | 环境阻塞数量 | blocked\_before\_runtime\_count | 74 | artifacts/phase12\_2\_clean/validation/verification/phase12\_summary.json | L4 | 环境阻塞样本不进入 runtime performance 分母。 |\par
| 第九章 | synthetic 样本数量 | synthetic\_sample\_count | 0 | artifacts/phase12\_2\_clean/validation/verification/phase12\_summary.json | L4 | clean validation 不包含 synthetic pipeline sample。 |\par
| 第九章 | 真实硬件未接触 | real\_controller\_contacted | False | artifacts/phase12\_2\_clean/validation/verification/phase12\_summary.json | L4 | 不能推出真实机械臂验收完成。 |\par
| 第九章 | 物理运动未发生 | hardware\_motion\_observed | False | artifacts/phase12\_2\_clean/validation/verification/phase12\_summary.json | L4 | 不能推出真实机械臂运动实验完成。 |\par
| 第九章 | 跨后端配对尚未 accepted | usable\_authoritative\_pair\_count | 0 | artifacts/phase12\_2\_clean/validation/verification/phase12\_summary.json | L4 | Isaac blocked 时不能声明 sim-to-sim 一致性。 |\par
| 第九章 | LLM-only fake provider 仅验证管线 | model\_runtime\_type | FAKE\_PROVIDER\_PIPELINE\_TEST | artifacts/thesis\_baselines/llm\_only/verification/llm\_only\_verification.json | L3 | fake provider 不代表真实大模型效果。 |\par

\section{附录 G 文献缺口}

正式参考文献只纳入已核验条目。当前需要继续补充云边协同机器人、快慢双系统、事件触发控制、LLM for robotics、LLM-only agent control、机器人安全、Sim2Real、MuJoCo、Isaac Sim、ROS 2、MoveIt、边缘自治和工程证据链相关文献。\par
